\documentclass[11pt]{article}

\usepackage{amsmath,amssymb,amsfonts}
\usepackage{bm}
\usepackage{geometry}
\usepackage{setspace}
\usepackage{titling}
\usepackage{titlesec}
\usepackage{microtype}
\usepackage{etoolbox}

% Page layout (matches technical note style)
\geometry{margin=1.15in}

% Section numbering depth
\setcounter{secnumdepth}{3}
\setcounter{tocdepth}{3}

% Compact title block
\setlength{\droptitle}{-3em}
\pretitle{\begin{center}\Large\bfseries}
\posttitle{\par\end{center}\vspace{-1em}}

% Line spacing similar to original PDF
\setstretch{1.08}

% More breathing room around equations
\setlength{\abovedisplayskip}{10pt}
\setlength{\belowdisplayskip}{10pt}

% Space around section titles
\titlespacing*{\section}{0pt}{14pt}{6pt}
\titlespacing*{\subsection}{0pt}{10pt}{4pt}

% Space around verbatim blocks
\AtBeginEnvironment{verbatim}{\vspace{6pt}}
\AtEndEnvironment{verbatim}{\vspace{6pt}}

\title{The ACDE Model with Trio and Sibling Data}
\date{}

\begin{document}
\maketitle
\vspace{-1em}
\tableofcontents
\vspace{1em}

\section{Introduction}

Classical twin studies partition phenotypic variance into additive genetic (A), 
shared environmental (C), and unique environmental (E) components. 
In many settings, however, dominance genetic effects (D) may also contribute 
substantially to phenotypic variation. 

The ACDE model extends the classical ACE framework by including dominance 
genetic variance. Estimation of all four variance components is not possible 
with twin data alone because C and D induce similar covariance patterns. 
Additional pedigree information is therefore required.

This document develops a compact formulation of the ACDE model using
trios (mother, father, child) and sibling data. The goal is to derive:

\begin{itemize}
\item the quantitative genetic model,
\item pedigree covariance expectations,
\item identifiability conditions,
\item a mixed model representation,
\item practical implementation and simulation.
\end{itemize}

\section{Quantitative Genetic Model}

For an individual $i$, the phenotype is written as
\begin{equation}
Y_i = \mu + A_i + D_i + C_i + E_i .
\end{equation}

The variance components are defined as
\begin{align}
A_i &\sim N(0,\sigma_A^2),\\
D_i &\sim N(0,\sigma_D^2),\\
C_i &\sim N(0,\sigma_C^2),\\
E_i &\sim N(0,\sigma_E^2).
\end{align}

The total phenotypic variance is therefore
\begin{equation}
\sigma_P^2 = \sigma_A^2 + \sigma_D^2 + \sigma_C^2 + \sigma_E^2 .
\end{equation}

Additive effects represent the sum of allele substitution effects across loci.  
Dominance effects represent within-locus interactions between alleles.  
Shared environment captures environmental influences common to members of 
the same household.  
Unique environment includes measurement error and individual-specific factors.

\section{Pedigree Covariances}

Covariances between relatives follow directly from kinship coefficients.  
Let $\phi_{ij}$ denote the kinship coefficient between individuals $i$ and $j$, 
and let $\Delta_{ij}$ denote the probability of sharing both alleles identical 
by descent at a locus.

\subsection{Additive covariance}

The additive genetic covariance between individuals $i$ and $j$ is
\begin{equation}
\mathrm{Cov}_A(i,j) = 2\phi_{ij}\sigma_A^2 .
\end{equation}

\subsection{Dominance covariance}

Dominance covariance depends on the probability of sharing two alleles 
identical by descent:
\begin{equation}
\mathrm{Cov}_D(i,j) = \Delta_{ij}\sigma_D^2 .
\end{equation}

\subsection{Shared environment}

Individuals raised in the same household share the same environmental effect:
\begin{equation}
\mathrm{Cov}_C(i,j) =
\begin{cases}
\sigma_C^2 & \text{if same household},\\
0 & \text{otherwise}.
\end{cases}
\end{equation}

\subsection{Key pedigree relationships}

For common family relationships we obtain:

\begin{center}
\begin{tabular}{lcc}
Relationship & $2\phi$ & $\Delta$\\
\hline
Parent--child & $1/2$ & $0$\\
Full siblings & $1/2$ & $1/4$\\
Unrelated & $0$ & $0$
\end{tabular}
\end{center}

These coefficients determine all expected covariances used in the ACDE model.

\section{Identifiability}

A central difficulty of the ACDE model is that shared environment (C) and 
dominance (D) produce similar covariance structures. Without sufficient 
pedigree information, the model parameters are not identifiable.

We therefore examine which family structures provide enough information to 
separate the four variance components.

\subsection{Trio data}

Consider a nuclear trio consisting of mother ($m$), father ($f$), and child ($c$).  
Assume parents are unrelated and do not share a household with the child 
after accounting for the child's shared environment.

The expected covariances are:

\begin{align}
\mathrm{Var}(Y_c) &= \sigma_A^2 + \sigma_D^2 + \sigma_C^2 + \sigma_E^2,\\
\mathrm{Var}(Y_m) &= \sigma_A^2 + \sigma_D^2 + \sigma_E^2,\\
\mathrm{Var}(Y_f) &= \sigma_A^2 + \sigma_D^2 + \sigma_E^2.
\end{align}

Parent--child covariance:
\begin{equation}
\mathrm{Cov}(Y_c,Y_m) = \frac{1}{2}\sigma_A^2.
\end{equation}

The same holds for the father.

\paragraph{Key observation}

Trio data identify the additive variance because parent–child covariance 
depends only on additive genetic effects. However, trio data alone cannot 
separate dominance and shared environmental variance because neither produces 
parent–child covariance.

\subsection{Trio + sibling data}

Now add a second child (siblings $c_1,c_2$).  
Sibling covariance becomes

\begin{equation}
\mathrm{Cov}(Y_{c_1},Y_{c_2})
= \frac{1}{2}\sigma_A^2 + \frac{1}{4}\sigma_D^2 + \sigma_C^2.
\end{equation}

This additional equation provides the crucial information required to 
separate dominance from shared environmental effects.

\subsection{Summary of identifying equations}

We now have three independent covariance relationships:

\begin{align}
\text{Parent--child:} \quad
& \frac{1}{2}\sigma_A^2, \\
\text{Sibling:} \quad
& \frac{1}{2}\sigma_A^2 + \frac{1}{4}\sigma_D^2 + \sigma_C^2, \\
\text{Individual variance:} \quad
& \sigma_A^2 + \sigma_D^2 + \sigma_C^2 + \sigma_E^2.
\end{align}

These equations are sufficient to identify all four variance components.

\section{Mixed Model Representation}

The ACDE model can be written as a linear mixed model.

\subsection{Vector notation}

Let $\mathbf{y}$ denote the vector of all phenotypes.  
The model is written as

\begin{equation}
\mathbf{y} = \mathbf{X}\bm{\beta}
+ \mathbf{Z_A}\mathbf{a}
+ \mathbf{Z_D}\mathbf{d}
+ \mathbf{Z_C}\mathbf{c}
+ \bm{\varepsilon}.
\end{equation}

Random effects are assumed to follow:

\begin{align}
\mathbf{a} &\sim N(0,\,\sigma_A^2\mathbf{A}),\\
\mathbf{d} &\sim N(0,\,\sigma_D^2\mathbf{D}),\\
\mathbf{c} &\sim N(0,\,\sigma_C^2\mathbf{I_f}),\\
\bm{\varepsilon} &\sim N(0,\,\sigma_E^2\mathbf{I}).
\end{align}

Here:

\begin{itemize}
\item $\mathbf{A}$ is the additive relationship matrix,
\item $\mathbf{D}$ is the dominance relationship matrix,
\item $\mathbf{I_f}$ indexes families.
\end{itemize}

\subsection{Phenotypic covariance matrix}

The marginal covariance of $\mathbf{y}$ becomes

\begin{equation}
\mathbf{V} =
\sigma_A^2\mathbf{A}
+ \sigma_D^2\mathbf{D}
+ \sigma_C^2\mathbf{C}
+ \sigma_E^2\mathbf{I}.
\end{equation}

This form enables estimation via REML using standard mixed-model software.

\section{Design Matrices}

To implement the ACDE mixed model, we must construct relationship matrices 
that encode pedigree structure.

\subsection{Additive relationship matrix}

The additive relationship matrix $\mathbf{A}$ contains twice the kinship 
coefficients. For the nuclear family with two siblings ordered as

\[
(m,\; f,\; c_1,\; c_2),
\]

the matrix is

\[
\mathbf{A} =
\begin{pmatrix}
1 & 0 & 1/2 & 1/2 \\
0 & 1 & 1/2 & 1/2 \\
1/2 & 1/2 & 1 & 1/2 \\
1/2 & 1/2 & 1/2 & 1
\end{pmatrix}.
\]

\subsection{Dominance relationship matrix}

Dominance relationships capture identity by descent of both alleles.  
For the same ordering, the dominance matrix becomes

\[
\mathbf{D} =
\begin{pmatrix}
1 & 0 & 0 & 0 \\
0 & 1 & 0 & 0 \\
0 & 0 & 1 & 1/4 \\
0 & 0 & 1/4 & 1
\end{pmatrix}.
\]

Parents do not share dominance covariance with their children.

\subsection{Shared environment matrix}

Individuals raised in the same household share a common environmental effect.  
For a family with two siblings,

\[
\mathbf{C} =
\begin{pmatrix}
0 & 0 & 0 & 0 \\
0 & 0 & 0 & 0 \\
0 & 0 & 1 & 1 \\
0 & 0 & 1 & 1
\end{pmatrix}.
\]

This matrix indexes the family-level random effect.

\section{R Implementation}

The ACDE model can be estimated using standard mixed model software such as 
\texttt{lme4} or \texttt{sommer}. The key requirement is the ability to supply 
custom covariance matrices.

\subsection{Unified model fitting function}

Below is a minimal implementation using the \texttt{sommer} package.

\begin{verbatim}
library(sommer)

fit_acde <- function(df, A, D, C) {

  fit <- mmer(
    y ~ 1,
    random = ~ vs(id, Gu = A)      # additive
            + vs(id, Gu = D)      # dominance
            + vs(fam)             # shared environment
            ,
    rcov = ~ units,
    data = df
  )

  return(fit)
}
\end{verbatim}

\subsection{Extracting variance components}

\begin{verbatim}
extract_vc <- function(fit) {
  vc <- summary(fit)$varcomp
  return(vc[,1])
}
\end{verbatim}

\subsection{Model comparison}

Nested models allow testing of variance components.

\begin{verbatim}
fit_ADE <- mmer(y ~ 1,
                random = ~ vs(id, Gu=A) + vs(id, Gu=D),
                rcov = ~ units, data=df)

anova(fit_ADE, fit_ACDE)
\end{verbatim}

Likelihood ratio tests provide formal evidence for the presence of shared 
environmental effects.

\section{Simulation Study}

A simulation study is useful to verify that the ACDE model can recover the 
true variance components under realistic sampling conditions.

\subsection{Simulation design}

We simulate nuclear families consisting of:
\begin{itemize}
\item mother,
\item father,
\item two siblings.
\end{itemize}

True variance components are set to

\[
\sigma_A^2 = 0.4,\quad
\sigma_D^2 = 0.2,\quad
\sigma_C^2 = 0.2,\quad
\sigma_E^2 = 0.2.
\]

These values produce substantial contributions from all sources of variation.

\subsection{Simulation algorithm}

For each family $j$:

\begin{enumerate}
\item Draw parental additive values
\[
A_m, A_f \sim N(0,\sigma_A^2).
\]

\item Generate offspring additive values using Mendelian sampling:
\[
A_{c_k} = \frac{1}{2}(A_m + A_f) + M_k,
\quad M_k \sim N\!\left(0,\frac{1}{2}\sigma_A^2\right).
\]

\item Simulate dominance deviations:
\[
D_i \sim N(0,\sigma_D^2).
\]

\item Draw family shared environment:
\[
C_j \sim N(0,\sigma_C^2).
\]

\item Draw individual environmental noise:
\[
E_i \sim N(0,\sigma_E^2).
\]

\item Construct phenotypes:
\[
Y_i = A_i + D_i + C_j + E_i.
\]
\end{enumerate}

\subsection{R simulation code}

\begin{verbatim}
simulate_family <- function() {

  Am <- rnorm(1,0,sqrt(.4))
  Af <- rnorm(1,0,sqrt(.4))

  Cfam <- rnorm(1,0,sqrt(.2))

  make_child <- function(){
    A <- 0.5*(Am+Af) + rnorm(1,0,sqrt(.2))
    D <- rnorm(1,0,sqrt(.2))
    E <- rnorm(1,0,sqrt(.2))
    return(A + D + Cfam + E)
  }

  Ym <- Am + rnorm(1,0,sqrt(.2)) + rnorm(1,0,sqrt(.2))
  Yf <- Af + rnorm(1,0,sqrt(.2)) + rnorm(1,0,sqrt(.2))

  Yc1 <- make_child()
  Yc2 <- make_child()

  return(c(Ym,Yf,Yc1,Yc2))
}
\end{verbatim}

\subsection{Model recovery experiment}

We simulate many families and fit the ACDE model.

\begin{verbatim}
Y <- replicate(1000, simulate_family())
\end{verbatim}

Estimated variance components should converge to the true values as the 
sample size increases.

\subsection{Expected results}

With sufficient sample size, REML estimates are approximately unbiased:
\[
\hat{\sigma}_A^2 \approx 0.4,\quad
\hat{\sigma}_D^2 \approx 0.2,\quad
\hat{\sigma}_C^2 \approx 0.2,\quad
\hat{\sigma}_E^2 \approx 0.2.
\]

\section{Interpretation}

The ACDE framework allows separation of additive, dominance, and shared 
environmental effects using standard family data.

Key insights include:

\begin{itemize}
\item Parent--child covariance identifies additive genetic variance.
\item Sibling covariance separates dominance from shared environment.
\item Mixed models provide a convenient estimation framework.
\end{itemize}

The approach extends classical twin models while retaining a simple 
likelihood-based implementation.

\section{Conclusion}

Trio and sibling data provide sufficient information to estimate the full 
ACDE variance decomposition. The mixed model formulation allows practical 
implementation using existing software and enables flexible extensions 
to larger pedigrees and genomic relationship matrices.

\end{document}
